% Options for packages loaded elsewhere
\PassOptionsToPackage{unicode}{hyperref}
\PassOptionsToPackage{hyphens}{url}
\documentclass[
  ignorenonframetext,
]{beamer}
\newif\ifbibliography
\usepackage{pgfpages}
\setbeamertemplate{caption}[numbered]
\setbeamertemplate{caption label separator}{: }
\setbeamercolor{caption name}{fg=normal text.fg}
\beamertemplatenavigationsymbolsempty
% remove section numbering
\setbeamertemplate{part page}{
  \centering
  \begin{beamercolorbox}[sep=16pt,center]{part title}
    \usebeamerfont{part title}\insertpart\par
  \end{beamercolorbox}
}
\setbeamertemplate{section page}{
  \centering
  \begin{beamercolorbox}[sep=12pt,center]{section title}
    \usebeamerfont{section title}\insertsection\par
  \end{beamercolorbox}
}
\setbeamertemplate{subsection page}{
  \centering
  \begin{beamercolorbox}[sep=8pt,center]{subsection title}
    \usebeamerfont{subsection title}\insertsubsection\par
  \end{beamercolorbox}
}
% Prevent slide breaks in the middle of a paragraph
\widowpenalties 1 10000
\raggedbottom
\AtBeginPart{
  \frame{\partpage}
}
\AtBeginSection{
  \ifbibliography
  \else
    \frame{\sectionpage}
  \fi
}
\AtBeginSubsection{
  \frame{\subsectionpage}
}
\usepackage{iftex}
\ifPDFTeX
  \usepackage[T1]{fontenc}
  \usepackage[utf8]{inputenc}
  \usepackage{textcomp} % provide euro and other symbols
\else % if luatex or xetex
  \usepackage{unicode-math} % this also loads fontspec
  \defaultfontfeatures{Scale=MatchLowercase}
  \defaultfontfeatures[\rmfamily]{Ligatures=TeX,Scale=1}
\fi
\usepackage{lmodern}
\ifPDFTeX\else
  % xetex/luatex font selection
\fi
% Use upquote if available, for straight quotes in verbatim environments
\IfFileExists{upquote.sty}{\usepackage{upquote}}{}
\IfFileExists{microtype.sty}{% use microtype if available
  \usepackage[]{microtype}
  \UseMicrotypeSet[protrusion]{basicmath} % disable protrusion for tt fonts
}{}
\makeatletter
\@ifundefined{KOMAClassName}{% if non-KOMA class
  \IfFileExists{parskip.sty}{%
    \usepackage{parskip}
  }{% else
    \setlength{\parindent}{0pt}
    \setlength{\parskip}{6pt plus 2pt minus 1pt}}
}{% if KOMA class
  \KOMAoptions{parskip=half}}
\makeatother
\usepackage{longtable,booktabs,array}
\usepackage{caption}
\captionsetup[table]{skip=6pt}
\newcounter{none} % for unnumbered tables
\usepackage{calc} % for calculating minipage widths
\usepackage{caption}
% Make caption package work with longtable
\makeatletter
\def\fnum@table{\tablename~\thetable}
\makeatother
\setlength{\emergencystretch}{3em} % prevent overfull lines
\providecommand{\tightlist}{%
  \setlength{\itemsep}{0pt}\setlength{\parskip}{0pt}}
% Auto-generated by infrastructure/rendering/_pdf_latex_helpers.py::extract_math_font_preamble.
% Minimal header passed to Pandoc via -H so Beamer slides pick up the same math font
% as the combined-PDF path. Adjust the manuscript's preamble.md to change the font globally.
\usepackage{unicode-math}
\setmathfont{latinmodern-math.otf}

% Slide layout defaults for warning-clean scientific decks.
\usepackage{etoolbox}
\IfFileExists{xurl.sty}{\usepackage{xurl}}{}
\IfFileExists{seqsplit.sty}{\usepackage{seqsplit}}{\newcommand{\seqsplit}[1]{#1}}
\protected\def\breakseq#1{\seqsplit{#1}}
\protected\def\breaktt#1{\begingroup\ttfamily\seqsplit{#1}\endgroup}
\setlength{\emergencystretch}{6em}
\tolerance=5000
\hbadness=10000
\hfuzz=1pt
\setlength{\tabcolsep}{2pt}
\AtBeginEnvironment{longtable}{\tiny\renewcommand{\arraystretch}{0.86}\setlength{\tabcolsep}{1pt}}
\AtBeginEnvironment{tabular}{\tiny\renewcommand{\arraystretch}{0.86}\setlength{\tabcolsep}{1pt}}
\AtBeginEnvironment{equation}{\tiny}
\AtBeginEnvironment{equation*}{\tiny}
\AtBeginEnvironment{align}{\tiny}
\AtBeginEnvironment{align*}{\tiny}
\AtBeginEnvironment{itemize}{\footnotesize}
\AtBeginEnvironment{enumerate}{\footnotesize}
\AtBeginEnvironment{description}{\footnotesize}
\setbeamerfont{caption}{size=\tiny}
\setbeamerfont{caption name}{size=\tiny}
\setbeamerfont{normal text}{size=\small}
\setbeamerfont{frametitle}{size=\small}
\setbeamerfont{section title}{size=\footnotesize}
\setbeamerfont{subsection title}{size=\footnotesize}
\setbeamertemplate{section page}{%
  \centering
  \begin{beamercolorbox}[sep=12pt,center,wd=\paperwidth]{section title}
    \parbox{0.86\paperwidth}{\centering\usebeamerfont{section title}\insertsection\par}
  \end{beamercolorbox}
}
\setbeamertemplate{subsection page}{%
  \centering
  \begin{beamercolorbox}[sep=8pt,center,wd=\paperwidth]{subsection title}
    \parbox{0.86\paperwidth}{\centering\usebeamerfont{subsection title}\insertsubsection\par}
  \end{beamercolorbox}
}
\setlength{\abovecaptionskip}{2pt}
\setlength{\belowcaptionskip}{0pt}

% Natbib and cross-reference fallbacks — slides don't load natbib
% or cleveref, but combined-PDF manuscript prose may emit these
% commands. The fallback renders citations as a bracketed key list
% and cross-references as detokenized labels so slides don't fail on
% undefined control sequences or raw underscores. \providecommand is
% a no-op if packages load later.
\providecommand{\citep}[1]{[#1]}
\providecommand{\citet}[1]{#1}
\providecommand{\citealp}[1]{#1}
\providecommand{\citeauthor}[1]{#1}
\providecommand{\citeyear}[1]{#1}
\providecommand{\cref}[1]{\texttt{\detokenize{#1}}}
\providecommand{\Cref}[1]{\texttt{\detokenize{#1}}}

\newtheorem{proposition}{Proposition}
\newtheorem{hypothesis}{Hypothesis}
\usepackage{bookmark}
\IfFileExists{xurl.sty}{\usepackage{xurl}}{} % add URL line breaks if available
\urlstyle{same}
% fallback for those not using the hyperref driver hyperxmp:
\makeatletter
\@ifundefined{xmpquote}{\newcommand{\xmpquote}[1]{#1}}{}
\makeatother
\hypersetup{
  hidelinks,
  pdfcreator={LaTeX via pandoc}}

\author{\texorpdfstring{}{}}
\date{}

\begin{document}

\section{Reproducibility and Release Evidence: Regenerating the Artifact
Chain}\label{sec:reproducibility}

\begin{frame}[fragile,allowframebreaks]{Reproducibility and Release
Evidence: Regenerating the Artifact Chain}
Every artifact in this work is regeneratable from source with fixed
seeds; nothing in \texttt{output/} is authoritative, and nothing
learner-facing is hand-typed.
\end{frame}

\subsection{Deterministic kernel outputs, stable graph ordering, and
rerunnable
counts}\label{deterministic-kernel-outputs-stable-graph-ordering-and-rerunnable-counts}

\begin{frame}[fragile,allowframebreaks]{Deterministic kernel outputs,
stable graph ordering, and rerunnable counts}
The kernels use natural-log units and either no randomness or an
explicitly seeded NumPy generator; the worked examples register rounded
values so two runs of \breaktt{build\_claim\_registry()} are
byte-identical. The curriculum graph's topological order processes ready
nodes in sorted-id order, so the exported SkillTree JSON is stable
across runs.
\end{frame}

\subsection{Verification gates for code, prose, artifacts, and rendered
surfaces}\label{verification-gates-for-code-prose-artifacts-and-rendered-surfaces}

\begin{frame}[fragile,allowframebreaks]{Verification gates for code,
prose, artifacts, and rendered surfaces}
The gate tables treat the curriculum build as scientific computing:
reported values should be traceable to workflow steps, custom code
should be tested and versioned, and research claims should be reported
with enough transparency that readers can distinguish reproducible local
artifacts from future replication or learner-outcome evidence (Sandve et
al. 2013; Wilson et al. 2014; Munafò et al. 2017).

The gates split into two families. The first runs anywhere the
repository can be checked out, with no Docker daemon and no browser.

\end{frame}

\begin{frame}[fragile,allowframebreaks]{Verification gates for code,
prose, artifacts, and rendered surfaces}
{\def\LTcaptype{none} % do not increment counter
\begin{longtable}[]{@{}
  >{\raggedright\arraybackslash}p{(\linewidth - 4\tabcolsep) * \real{0.3333}}
  >{\raggedright\arraybackslash}p{(\linewidth - 4\tabcolsep) * \real{0.3333}}
  >{\raggedright\arraybackslash}p{(\linewidth - 4\tabcolsep) * \real{0.3333}}@{}}
\toprule\noalign{}
\begin{minipage}[b]{\linewidth}\raggedright
Offline gate
\end{minipage} & \begin{minipage}[b]{\linewidth}\raggedright
Command
\end{minipage} & \begin{minipage}[b]{\linewidth}\raggedright
Evidence
\end{minipage} \\
\midrule\noalign{}
\endhead
Tests + coverage &
\breaktt{uv\ run\ pytest\ -\/-cov=src\ -\/-cov-fail-under=90} & enforced
aggregate floor: 90\%; use command output as the current coverage
number, and read the module caveat below \\
Kernel coverage & within the coverage run & kernel package coverage is
checked inside the same suite rather than asserted from prose \\
Lint & \breaktt{uvx\ ruff\ check\ src\ scripts\ tests} & clean
required \\
Format & \breaktt{uvx\ ruff\ format\ -\/-check\ src\ scripts\ tests} &
clean required \\
Types & \breaktt{uv\ run\ mypy\ src} & no issues required \\
No mocks &
\texttt{rg\ "unittest\textbar{}MagicMock\textbar{}@patch\textbar{}autospec"\ tests/}
& no mock factories should be present \\
Provenance & \breaktt{scripts/run\_kernels.py} and
\breaktt{scripts/export\_skilltree.py} & every claim token resolves and
the bare-number linter reports zero violations \\
Scholarship & every authored citation maps to a BibTeX entry and
source-role lane & enforced before figures \\
Figure artifacts & generated figure registry, manuscript
labels/cross-references, PNG information content, captions, alt text,
and dashboard figure artifacts agree & enforced by
\texttt{run\_kernels.py} and
\breaktt{validate\_manuscript\_artifacts.py} \\
Image pinning &
\breaktt{uv\ run\ python\ scripts/validate\_docker\_digests.py} & every
Compose image reference carries a \texttt{sha256} digest; the file is
parsed, no daemon is contacted \\
Release inventory &
\breaktt{uv\ run\ python\ scripts/validate\_release\_readiness\_inventory.py}
& checked against live git status, untracked files, ignored outputs and
blocker probes \\
Manuscript render & parent template renderer & citation and figure
references must resolve \\
\bottomrule\noalign{}
\end{longtable}
}

\end{frame}

\begin{frame}[fragile,allowframebreaks]{Verification gates for code,
prose, artifacts, and rendered surfaces}
The second family is live. Every row below needs a reachable Docker
daemon, and the browser rows additionally need the
\breaktt{chrome-devtools-axi} Chrome driver. These gates fail closed
rather than downgrading to the artifact-only dashboard, which is why
they are listed separately instead of folded into the table above.

\end{frame}

\begin{frame}[fragile,allowframebreaks]{Verification gates for code,
prose, artifacts, and rendered surfaces}
{\def\LTcaptype{none} % do not increment counter
\begin{longtable}[]{@{}
  >{\raggedright\arraybackslash}p{(\linewidth - 4\tabcolsep) * \real{0.3333}}
  >{\raggedright\arraybackslash}p{(\linewidth - 4\tabcolsep) * \real{0.3333}}
  >{\raggedright\arraybackslash}p{(\linewidth - 4\tabcolsep) * \real{0.3333}}@{}}
\toprule\noalign{}
\begin{minipage}[b]{\linewidth}\raggedright
Live gate
\end{minipage} & \begin{minipage}[b]{\linewidth}\raggedright
Command
\end{minipage} & \begin{minipage}[b]{\linewidth}\raggedright
Evidence
\end{minipage} \\
\midrule\noalign{}
\endhead
Live dependency parity &
\breaktt{uv\ run\ python\ scripts/reconcile\_live\_dependencies.py\ -\/-check-only}
& live prerequisite edges compared with the exported project JSON
without writing edges \\
Live dashboard wiring &
\breaktt{uv\ run\ python\ scripts/validate\_live\_skilltree\_dashboard.py\ -\/-require-live-admin\ -\/-require-live-dependencies\ -\/-require-figure-contracts}
& live admin skill-count parity, exact live dependency parity, and
figure-contract agreement \\
Raw SkillTree login &
\breaktt{uv\ run\ python\ scripts/validate\_skilltree\_form\_login.py} &
a real browser drives the \texttt{FORM} login through to the
administrator surface; the JSON summary is redacted \\
Dashboard console &
\breaktt{uv\ run\ python\ scripts/validate\_dashboard\_console.py\ -\/-require-browser\ -\/-require-dashboard}
& no console errors except the allow-listed upstream \texttt{Inception}
status probe \\
App-ready stack &
\breaktt{uv\ run\ python\ scripts/validate\_full\_app\_stack.py\ -\/-app-ready\ -\/-skip-regenerate}
& Compose bring-up, seed, dashboard, login, and console pass; the JSON
summary is labeled app-ready \\
Full-parity stack &
\breaktt{uv\ run\ python\ scripts/validate\_full\_app\_stack.py\ -\/-skip-regenerate}
& the same sequence plus complete live dependency reconciliation; the
JSON summary is labeled full-parity \\
\bottomrule\noalign{}
\end{longtable}
}
\end{frame}

\begin{frame}[fragile,allowframebreaks]{Coverage floor caveat}
\protect\phantomsection\label{coverage-floor-caveat}
The coverage floor is enforced on the aggregate, and the aggregate is
not the same claim as every module meeting it. Four modules currently
sit below the floor while the suite as a whole clears it:
\breaktt{src/integration/full\_app\_stack.py},
\breaktt{src/skilltree/form\_login\_validation.py},
\breaktt{src/scholarship/source\_health\_network.py}, and
\breaktt{src/dashboard/payload.py}. Each one is a boundary to something
the offline suite cannot reach --- a Docker daemon, a browser
subprocess, live HTTP, or the local-auth server path --- and each is
exercised instead by the live gates in the second table. Two defects
show why that exception is a routing decision rather than an excuse: a
dashboard content-security-policy directive that blanked the official
SkillTree panel was invisible to every offline test and surfaced only in
the browser console gate, and a reconciler abort-on-slow-write surfaced
\end{frame}

\begin{frame}[fragile,allowframebreaks]{Coverage floor caveat}
only during a complete live dependency run. Read the aggregate number as
a property of the suite, not as evidence of uniform module coverage.
\end{frame}

\subsection{Release-evidence ledger and publication
blockers}\label{release-evidence-ledger-and-publication-blockers}

\begin{frame}[fragile,allowframebreaks]{Release inventory comparison}
\protect\phantomsection\label{release-inventory-comparison}
Release evidence is recorded as a local, machine-checked ledger rather
than as a clean public-release assertion. During review, the checkout
may intentionally contain dirty or untracked files;
\breaktt{docs/release\_readiness\_inventory.md} is the authoritative
inventory for that state and is regenerated with
\breaktt{uv\ run\ python\ scripts/validate\_release\_readiness\_inventory.py\ -\/-write}.

\end{frame}

\begin{frame}[fragile,allowframebreaks]{Release inventory comparison}
The default validator then compares the document with live
\breaktt{git\ status\ -\/-short},
\breaktt{git\ ls-files\ -o\ -\/-exclude-standard},
\breaktt{git\ diff\ -\/-stat}, ignored generated-output examples, and the
current external blocker probes. Those probes report three distinct
positions, and collapsing them into one blocked-or-ready summary would
misstate all three.

\texttt{docker\ info} reaches the Docker/Colima daemon, so the
Docker-dependent live rehearsal is no longer blocked. It has been run:
\breaktt{validate\_full\_app\_stack.py\ -\/-skip-regenerate} completed
against the pinned SkillTree image in full-parity mode, with the live
dependency graph matching the export edge for edge, the live admin skill
count matching the exported catalog, every registered artifact and
figure contract agreeing, the raw \texttt{FORM} login reaching the
administrator surface, and the browser console gate clean. What remains
\end{frame}

\begin{frame}[fragile,allowframebreaks]{Release inventory comparison}
is to rerun that sequence as a release handoff rather than as a
development rehearsal, on the commit that is actually released.

\texttt{git\ ls-remote} against the public
\breaktt{ActiveInferenceInstitute/Active\_Skillference} target now
succeeds, and that is precisely why the validator distinguishes
reachability from publication: the target exists and holds no refs. An
exit-code-only check would read that empty placeholder as ready. The
inventory therefore reports it as reachable but empty, and DOI,
archive-deposit, and public-repository metadata fields stay blank until
refs actually land there. Nothing in this repository is published at
that address today.

\end{frame}

\begin{frame}[fragile,allowframebreaks]{Release inventory comparison}
The independent cross-vendor audit is closed in \texttt{tasks.yaml}, not
pending. Its recorded scope is a point-in-time, read-only static review
of the dashboard, local-auth, and full-stack sources and their tests,
and within that scope it returned no unresolved critical or high
findings. It did not cover dependency and supply-chain exposure, secrets
in git history, or dynamic runtime testing, so it closes a specific
review task rather than establishing operational assurance. Branch
hygiene remains open while the checkout still carries modified tracked
paths.

\end{frame}

\begin{frame}[fragile,allowframebreaks]{Release inventory comparison}
Reproducible-workflow guidance motivates the rerunnable command ledger,
versioned code and transparent artifact state (Sandve et al. 2013;
Wilson et al. 2014), while FAIR and transparency-oriented open-science
standards motivate the metadata and evidence-ledger shape (Wilkinson et
al. 2016; Nosek et al. 2015; Munafò et al. 2017). These sources do not
turn a dirty local checkout into a public archive. Software-citation and
machine-readable metadata guidance define the repository-level citation
files and release metadata surfaces (Smith et al. 2016; Citation File
Format Project 2026; CodeMeta Project 2026), while GitHub and Zenodo
documentation keep DOI language tied to a real public target and deposit
(GitHub Docs 2026; Zenodo 2026). SPDX, OSI and REUSE references bound
the MIT and third-party-notice signposting (SPDX Workgroup 2026; Open
Source Initiative 2026; Free Software Foundation Europe 2024); they do
\end{frame}

\begin{frame}[fragile,allowframebreaks]{Release inventory comparison}
not provide legal advice, relicense third-party materials, or remove the
need for separate deployment review. A release record should include the
base commit hash, whether the worktree was dirty, the generated
inventory, the commands in the gate table above, and the artifact
manifests in \texttt{output/data/} and
\breaktt{output/figures/figure\_registry.json}; it should not promote
this local evidence into a DOI or public-repository claim until the
external blockers are resolved.

\end{frame}

\begin{frame}[fragile,allowframebreaks]{Release inventory comparison}
The release inventory also carries a machine-checked \textbf{Key
Artifact Hashes} section. During validation, it records SHA-256 digests
and byte sizes for \breaktt{output/data/claims.json},
\breaktt{output/data/manuscript\_variables.json},
\breaktt{output/data/skilltree\_project.json},
\breaktt{output/figures/figure\_registry.json}, and every registered
figure PNG currently named by the figure registry. The same validator
rejects stale hash rows, so a regenerated artifact cannot silently drift
from the manuscript's local reproducibility evidence. This is still
local evidence: without a release tag, public archive, and clean
external blocker probes, it remains a checkout ledger rather than an
immutable publication bundle.
\end{frame}

\subsection{Provenance chain from kernel function to learner-facing
value}\label{provenance-chain-from-kernel-function-to-learner-facing-value}

\begin{frame}[fragile,allowframebreaks]{Provenance chain from kernel
function to learner-facing value}
The chain from kernel to learner is auditable end to end: a kernel
function computes a value -\textgreater{} \texttt{examples.py} registers
it as a named \texttt{Claim} with its source function -\textgreater{} a
skill or manuscript section references it as \texttt{claim:NAME}
-\textgreater{} the exporter or render hook resolves it, failing closed
if it cannot -\textgreater{} the value appears in the SkillTree skill
description, the manuscript, and \breaktt{output/data/claims.json}.
Changing a kernel changes the claim, the rendered content, and the
figures together. The separate manuscript artifact audit closes the
non-numeric drift path: a figure cannot be registered without a
manuscript declaration, referenced without a declaration, missing from
disk, or absent from the dashboard figure inventory.
\end{frame}

\subsection{Disposable generated artifacts and output
inventory}\label{disposable-generated-artifacts-and-output-inventory}

\begin{frame}[fragile,allowframebreaks]{Disposable generated artifacts
and output inventory}
\breaktt{scripts/run\_kernels.py} writes \breaktt{output/data/claims.json}
(the provenance ledger), \breaktt{output/data/manuscript\_variables.json}
(claims plus curriculum counts),
\breaktt{output/data/scholarship\_coverage.json} (the citation lane
audit), \breaktt{output/data/scholarship\_gap\_report.json} (the
scholarship frontier-source and claim-boundary audit),
\breaktt{output/data/manuscript\_artifacts.json} (the figure-surface
audit), \breaktt{output/data/cover\_asset\_audit.json} (the first-page
cover audit), \breaktt{output/data/front\_matter\_visual\_audit.json}
(the page-two atlas audit),
\breaktt{output/data/course\_material\_index.json} (local modular lesson
records), and \breaktt{output/data/course\_material\_audit.json} (the
lesson-record completeness audit), plus 21 registered figures plus the
generated cover and named-skill atlas PNGs under
\texttt{output/figures/}. \breaktt{scripts/export\_skilltree.py} writes
\end{frame}

\begin{frame}[fragile,allowframebreaks]{Disposable generated artifacts
and output inventory}
\breaktt{output/data/skilltree\_project.json}, the import-ready SkillTree
project. \breaktt{scripts/z\_generate\_manuscript\_variables.py} is the
parent-template render hook: it runs the kernel/provenance gate and
writes resolved manuscript copies under \breaktt{output/manuscript/}
before PDF, HTML, slide, DOCX, or EPUB rendering.
\end{frame}

\end{document}
